% ═══════════════════════════════════════════════════════════════════
%  Conclusion générale et perspectives
% ═══════════════════════════════════════════════════════════════════
\pageLiminaire{Conclusion Générale et Perspectives}

\section*{Bilan du travail réalisé}

Ce projet de fin d'études a abouti à une plate-forme logicielle complète pour
\yani\ : une API REST portant l'intégralité des règles métier, une application
mobile trilingue destinée aux clientes, et un back-office web destiné à
l'institut. Les trois applications partagent une même base PostgreSQL et un
même langage, et sont validées ensemble par une chaîne d'intégration continue.

Les objectifs fixés au départ ont été atteints. La réservation est possible à
toute heure, sur une grille de créneaux calculée à partir des horaires réels du
centre. Le carnet papier est remplacé par une base sauvegardée, consultable
depuis plusieurs postes. La carte de fidélité tamponnée a laissé place à un
compte doté d'un solde, d'un historique, de paliers de visites et de bons
traçables. Le stock est suivi, contrôlé et décrémenté au bon moment. La gérante
dispose enfin d'une visibilité sur son activité et de l'export de ses données.

Le travail ne s'est cependant pas limité à la mise en \oe uvre de
fonctionnalités. Sa part la plus exigeante a porté sur des propriétés qui ne se
voient pas à l'écran :

\begin{itemize}[leftmargin=1.2cm]
  \item \textbf{la correction sous concurrence} --- un créneau n'est jamais
  vendu deux fois, une unité de stock jamais attribuée deux fois, un point de
  fidélité jamais crédité deux fois ; ces propriétés sont établies par des
  tests qui écrivent dans une véritable base PostgreSQL, parce qu'aucune base
  simulée ne saurait les démontrer ;

  \item \textbf{la fidélité au temps réel} --- les instants sont stockés en
  temps universel, les horaires exprimés en heure locale du centre, et la
  conversion n'est faite qu'à un seul endroit, par le serveur, qui possède la
  base des fuseaux ;

  \item \textbf{la conservation de l'historique} --- une cliente peut exercer
  son droit à l'effacement sans que le chiffre d'affaires qu'elle a généré ne
  disparaisse avec elle ; les modes de suppression déclarés dans le schéma
  rendent cette garantie structurelle plutôt que déclarative ;

  \item \textbf{la sécurité} --- traitée à partir d'un modèle de menaces
  explicite et non comme une couche ajoutée en fin de parcours : hachage
  Argon2id, jetons courts, rotation avec révocation de famille en cas de
  réutilisation détectée, uniformisation du temps de réponse contre
  l'énumération de comptes, autorisation vérifiée par le serveur sur chaque
  route et doublée d'un contrôle de propriété des ressources, validation
  stricte des entrées, reconnaissance des fichiers téléversés par leur
  signature binaire.
\end{itemize}

\section*{Apports personnels}

Sur le plan technique, ce projet a été l'occasion de mener seul, de bout en
bout, une réalisation à trois cibles : un serveur, une application web et une
application mobile. Il m'a permis d'approfondir des sujets que le cursus
aborde nécessairement de façon théorique --- la concurrence en base de données,
la gestion des fuseaux horaires, la sécurité applicative --- et de découvrir
qu'ils ne se comprennent réellement qu'au moment où un test les met en échec.

L'apport le plus net concerne la sécurité, et il tient moins à des techniques
qu'à un changement de regard. Appliquer une grille de menaces à un système que
l'on a soi-même écrit oblige à l'examiner du point de vue de quelqu'un qui
cherche à en abuser, et non de celui qui l'utilise comme prévu. C'est cet
exercice qui a fait apparaître ce qu'aucune relecture fonctionnelle ne
signalait : que masquer un bouton n'a jamais protégé la route correspondante,
qu'un temps de réponse suffit à trahir l'existence d'un compte, et qu'une
condition de course déclenchable à volonté n'est pas un défaut de correction
mais une vulnérabilité.

Le test de concurrence sur les réservations en est l'illustration la plus
nette. L'implémentation initiale --- compter puis insérer --- paraissait
évidemment correcte à la lecture. Il a fallu trois réservations réellement
simultanées pour montrer qu'elle ne l'était pas. C'est l'enseignement que je
retiens le plus fermement de ce projet : \emph{en matière de concurrence,
l'intuition ne vaut rien, et seul un test qui reproduit la simultanéité fait
autorité}.

Sur le plan méthodologique, la démarche agile a démontré son intérêt de façon
concrète. Plusieurs des règles les plus importantes du système --- qu'un
rendez-vous occupe un créneau fixe et non la durée de sa prestation, qu'une
réclamation de récompense doive laisser une trace au comptoir, qu'un produit
retiré du catalogue ne doive pas s'effacer de l'historique --- ne figuraient
dans aucune spécification initiale. Elles sont apparues en montrant des
versions fonctionnelles à la personne qui allait s'en servir.

Sur le plan personnel, travailler pour une utilisatrice unique, identifiée et
non informaticienne a été formateur. Un message d'erreur incompréhensible n'est
pas un détail lorsque la personne qui le reçoit n'a personne à qui le
transmettre. Cette contrainte a orienté un grand nombre de décisions, du
libellé des statuts dans les exports Excel jusqu'au choix d'afficher les
créneaux indisponibles plutôt que de les masquer.

\section*{Difficultés rencontrées}

Trois difficultés méritent d'être mentionnées, ainsi que la façon dont elles
ont été surmontées.

\paragraph{La concurrence sur les réservations.} Décrite plus haut, elle a
imposé de descendre au niveau du SGBD et de recourir à un verrou consultatif
PostgreSQL --- une primitive qui ne figure dans aucune abstraction d'ORM.

\paragraph{Les fuseaux horaires.} Les erreurs de fuseau sont silencieuses :
elles ne lèvent aucune exception et produisent des résultats plausibles mais
faux, décalés d'une heure. La convention a donc été fixée explicitement ---
stockage en temps universel, horaires en heure locale, conversion en un seul
endroit --- et documentée dans le code comme dans l'API.

\paragraph{Les pièges de déploiement invisibles.} Plusieurs réglages
d'exploitation ne cassent rien immédiatement, mais produisent des symptômes
déroutants une fois en service : le réglage du proxy de confiance, sans lequel
tout l'institut partage un seul compteur de requêtes ; la liste des origines
autorisées, dont l'oubli n'apparaît que dans la console du navigateur. Ces
pièges ont été documentés à l'endroit où ils se manifestent, avec leur symptôme
et leur cause, plutôt que dans une documentation séparée que personne ne
consulte au moment utile.

\section*{Perspectives}

La plate-forme livrée est fonctionnelle, mais plusieurs évolutions ont été
identifiées, classées ici par horizon.

\paragraph{À court terme --- finaliser la mise en service.}
\begin{itemize}[leftmargin=1.2cm,nosep]
  \item Acquisition d'un nom de domaine et d'un certificat HTTPS, avec le
  reverse proxy correspondant.
  \item Mise en place d'un serveur de messagerie définitif.
  \item Publication de l'application sur les magasins Android et iOS.
\end{itemize}

\paragraph{À moyen terme --- enrichir l'usage.}
\begin{itemize}[leftmargin=1.2cm,nosep]
  \item \textbf{Notifications \textit{push}} : rappel de rendez-vous la veille,
  avis de commande prête, alerte de palier de fidélité franchi. C'est
  l'évolution la plus demandée, et elle dépend d'une compilation native, donc
  du nom de domaine.
  \item \textbf{Paiement en ligne}, en complément du paiement à la remise.
  \item \textbf{Statistiques de pilotage} : chiffre d'affaires par période,
  prestations les plus demandées, taux de retour de la clientèle. Les données
  sont déjà en base ; il ne manque que leur restitution.
  \item \textbf{Gestion fine des ressources} : aujourd'hui, la capacité du
  centre est un nombre unique de places réservables. Distinguer les cabines et
  les esthéticiennes permettrait de proposer des créneaux différenciés selon la
  prestation.
\end{itemize}

\paragraph{À moyen terme --- durcir.} Les limites de sécurité inventoriées au
chapitre~4 appellent quatre chantiers, par ordre de priorité décroissante : un
\textbf{second facteur d'authentification} sur le compte de direction, qui donne
accès à l'ensemble du fichier clientes ; un \textbf{test d'intrusion externe},
seule façon de découvrir ce que la personne qui a écrit le code ne pense pas à
chercher ; une \textbf{supervision} qui alerte sur les séries d'échecs de
connexion plutôt que de se contenter de les journaliser ; et le report du
compteur de débit dans un magasin partagé, préalable à toute mise à l'échelle
horizontale.

\paragraph{À plus long terme.}
\begin{itemize}[leftmargin=1.2cm,nosep]
  \item \textbf{Fiche cliente} tenue par l'institut : historique des soins,
  allergies signalées, préférences --- un dossier de suivi qui prolongerait la
  relation personnelle qui fait la valeur du centre.
  \item \textbf{Gestion multi-établissements}, si l'institut venait à ouvrir un
  second point de vente.
  \item \textbf{Automatisation du marketing} : relance des clientes non revenues
  depuis un certain temps, offres d'anniversaire.
\end{itemize}

\vspace{0.5cm}

Au terme de ce travail, la plate-forme est prête à entrer en service. Elle ne
remplace pas la relation entre l'institut et ses clientes --- elle ne le
cherche pas --- mais elle décharge la gérante de ce qui l'en éloignait : le
téléphone qui sonne pendant un soin, le carnet à consulter, la carte de
fidélité à retrouver. C'est, en définitive, la mesure la plus juste de sa
réussite.
